\chapter{Introdução}

Teste de software é uma das atividades presentes no processo de desenvolvimento e manutenção de software. Conforme o \textit{Guide to the Software Engineering Body of Knowledge} (SWEBOK) \cite{bourque:etal:2014} define, teste de software consiste na verificação dinâmica de que um programa provê um comportamento esperado a partir de um conjunto finito de casos de teste, adequadamente selecionados a partir de um domínio de execução potencialmente infinito. Por muito tempo, teste de software foi visto sobretudo como uma etapa de verificação da qualidade do software produzido, e não como meio de construir um produto com qualidade \cite{murugesan:1994}. Nesse sentido, testes eram frequentemente realizados apenas ao final da etapa de codificação \cite{yang:etal:2011}, com o objetivo de detectar falhas e fornecer um indicativo de que o software funciona conforme o esperado \cite{bertolino:2007}. Com o passar do tempo, a percepção sobre a utilidade do teste evoluiu. A atividade deixou de ser vista como etapa posterior ao desenvolvimento e passou a estar presente durante todo o processo de desenvolvimento e manutenção de software \cite{bourque:etal:2014}. \citeonline{meszaros:2007} considera a atividade de teste parte intrínseca do processo de desenvolvimento de software. Assim, a atividade de teste de software passou a incorporar propósitos variados, como detectar falhas \cite{tiwari:etal:2013}, prevenir que erros sejam introduzidos \cite{beizer:1990}, prover um indicativo de que o software funciona \cite{bertolino:2007}, auxiliar na compreensão do projeto e do código do sistema \cite{freeman:etal:2009} e servir como meio para especificação \cite{alvestad:2007} e documentação \cite{haugset:2008} dos requisitos.

Nesse contexto, o código de teste deixou de ser um artefato auxiliar e passou a exigir os mesmos cuidados aplicados ao código de produção. Assim como o código da aplicação, o código de teste precisa ser mantido, compreendido e ajustado durante todo o projeto \cite{greiler:etal:2013:b}. Para \citeonline{meszaros:2007}, a atividade de teste não deveria aumentar o esforço total de desenvolvimento e manutenção do software, uma vez que o esforço gasto com essa atividade deve ser compensado pela melhoria da qualidade do software. Entretanto, manter o código de teste apresenta desafios, principalmente quando a codificação e a evolução dos testes são realizadas diretamente pelo desenvolvedor. Conforme o software evolui e a quantidade de código aumenta, torna-se mais difícil manter os testes com consistência e sem perder o controle sobre eles. Fenômeno semelhante ocorre no código de produção: sem refatoração contínua, cada nova funcionalidade tende a exigir mais esforço do que a anterior \cite{fowler:1999}. Com os testes, observa-se comportamento análogo. Novos testes passam a ser mais difíceis de adicionar e a duplicação tende a aumentar \cite{berner:etal:2005}. Isso pode levar à degradação da manutenibilidade do código de teste \cite{greiler:etal:2013:a}. Segundo \citeonline{emery:2009}, um motivo comum para o abandono de testes automatizados é o fato de os testes se tornarem frágeis e de alto custo de manutenção, o que pode tornar os desenvolvedores reticentes tanto em adicionar novos testes quanto em executar os existentes.

Evitar a duplicação no código de teste é uma das formas de melhorar a manutenibilidade. Isso porque mudanças na aplicação passam a ter impacto menor nos testes, e o desenvolvedor dispõe de menos código para manter e gerenciar \cite{berner:etal:2005,greiler:etal:2013:a}. De acordo com \citeonline{berner:etal:2005}, testes tendem a ser repetitivos e apresentam alto potencial de reuso; conforme o software evolui e o número de testes aumenta, é comum que a quantidade de código duplicado entre os testes também aumente. Diante disso, técnicas que reduzam essa duplicação podem contribuir para preservar a manutenibilidade do código de teste ao longo do projeto.
